\documentclass[a4paper]{arrowhead}

\usepackage[yyyymmdd]{datetime}
\usepackage{etoolbox}
\usepackage[utf8]{inputenc}
\usepackage[T1]{fontenc}
\usepackage{multirow}
\usepackage{longtable}
\usepackage{tikz}
\usepackage{adjustbox}
\AtBeginEnvironment{longtable}{\small}
\BeforeBeginEnvironment{tikzpicture}{\begin{adjustbox}{max width=\textwidth}}
\AfterEndEnvironment{tikzpicture}{\end{adjustbox}}
\usetikzlibrary{positioning,shapes.multipart,arrows.meta}

\renewcommand{\dateseparator}{-}
\setlength{\emergencystretch}{3em}
\lstset{language=HTTP,breakatwhitespace=false,columns=fullflexible}

\newcommand{\fparam}[1]{\textit{\textcolor{ArrowheadBlue}{#1}}}

%% Special references
\newcommand{\mref}[2][]{{\textcolor{ArrowheadPurple}{\hyperref[sec:model:#2]{\ifstrempty{#1}{#2}{#1}}}}}
\newcommand{\oref}[1]{{\textcolor{ArrowheadBlue}{\hyperref[sec:operations:#1]{#1}}}}
\newcommand{\pdef}[1]{{\textcolor{ArrowheadGrey}{#1 \label{sec:model:primitives:#1} \label{sec:model:primitives:#1s}}}}
\newcommand{\pref}[1]{{\textcolor{ArrowheadGrey}{\hyperref[sec:model:primitives:#1]{#1}}}}

% as template_idd.tex; the input and output arguments are passed already formatted (\mref or \pref)
\newrobustcmd\osubsection[5]{
  \addtocounter{subsection}{1}
  \addcontentsline{toc}{subsection}{\protect\numberline{\thesubsection}operation \textcolor{ArrowheadBlue}{#1}}
  \subsection*{
    \thesubsection\quad
    #2 \textcolor{ArrowheadPurple}{#3} \\
    \small
    \hspace*{0.075\textwidth}\begin{minipage}{0.1\textwidth}
      \vspace*{1mm}
      Operation: \\
      \notblank{#4}{Input: \\}{}
      \notblank{#5}{Output: \\}{}
    \end{minipage}%
    \begin{minipage}{0.825\textwidth}
      \vspace*{1mm}
      \textcolor{ArrowheadBlue}{#1} \\
      \notblank{#4}{#4 \\}{}
      \notblank{#5}{#5 \\}{}
    \end{minipage}
  }
  \label{sec:operations:#1}
}
\newrobustcmd\msubsection[2]{
  \addtocounter{subsection}{1}
  \addcontentsline{toc}{subsection}{\protect\numberline{\thesubsection}#1 \textcolor{ArrowheadPurple}{#2}}
  \subsection*{\thesubsection\quad#1 \textcolor{ArrowheadPurple}{#2}}
  \label{sec:model:#2}
}
%%

\begin{document}
\sloppy

%% Arrowhead Document Properties
\ArrowheadTitle{authorizationToken HTTP/TLS/JSON}
\ArrowheadServiceID{authorizationToken}
\ArrowheadType{Interface Design Description}
\ArrowheadTypeShort{IDD}
\ArrowheadVersion{5.2.2}
\ArrowheadDate{2026-10-22}
\ArrowheadAuthor{Jerker Delsing}
\ArrowheadStatus{RELEASE CANDIDATE}
\ArrowheadContact{jerker.delsing@ltu.se}
\ArrowheadFooter{\href{www.arrowhead.eu}{www.arrowhead.eu}}
\ArrowheadSetup
%%

%% Front Page
\begin{center}
  \vspace*{1cm}
  \huge{\arrowtitle}

  \vspace*{0.2cm}
  \LARGE{\arrowtype}
  \vspace*{1cm}

  \Large{Service ID: \textit{"\arrowid"}}

  \vspace*{\fill}

  % Front Page Image

  \vspace*{1cm}
  \vspace*{\fill}

  % Front Page Abstract
  \begin{abstract}
    This document describes the HTTP/1.1 protocol with TLS payload security and JSON payload encoding variant of the \texttt{authorizationToken} service of the ConsumerAuthorization system.
  \end{abstract}

  \vspace*{1cm}
\end{center}

\newpage
%%

%% Table of Contents
\tableofcontents
\newpage
%%

\section{Overview}
\label{sec:overview}

This document describes the \texttt{authorizationToken} service interface of the ConsumerAuthorization system, which provides authorization tokens: a consumer obtains a token for a specific provider, service and operation; the provider verifies it.

It's implemented using protocol, encoding as stated in below table:
\begin{table}[ht!]
  \centering
  \begin{tabular}{|l|l|l|}
    \hline
    \rowcolor{gray!33} Profile type & Type & Version \\ \hline
    Transfer protocol & HTTP & 1.1 \cite{rfc9112} \\ \hline
    Data encryption & TLS & 1.3 \cite{rfc8446}, 1.2 \cite{rfc5246} \\ \hline
    Encoding & JSON & RFC 8259 \cite{rfc8259} \\ \hline
    Compression & N/A & - \\ \hline
    Semantics & N/A & - \\ \hline
    Ontology & N/A & - \\ \hline
  \end{tabular}
  \caption{Communication and semantics details used for the \texttt{authorizationToken} service interface}
  \label{tab:comunication_semantics_profile}
\end{table}

This document provides the Interface Design Description IDD to the \textit{authorizationToken -- Service Description} document \cite{ca-authorizationToken-sd}.
For further details about how this service is meant to be used, please consult that document.

The rest of this document describes how to realize the \texttt{authorizationToken} service
interface in detail. Both in terms of its operations (Section \ref{sec:operations}) and its information model (Section \ref{sec:model}).

\subsection{Interface registration}

The service is registered in the ServiceRegistry with the interface template \texttt{generic\_\allowbreak{}https} (\texttt{generic\_\allowbreak{}http} when TLS is disabled), access port 8445 and base path \texttt{/\allowbreak{}consumerauthorization/\allowbreak{}authorization-\allowbreak{}token}. The \texttt{operations} property of the interface lists the operations below with their method and path relative to the base path (path parameters are placeholders such as \texttt{\{id\}} that the consumer fills in):

\noindent\begin{longtable}{| p{5.2cm} | p{1.6cm} | p{8.4cm} |} \hline
\rowcolor{gray!33} Operation & Method & Path \\ \hline \endhead
\texttt{generate} & POST & \texttt{/\allowbreak{}generate} \\ \hline
\texttt{verify} & GET & \texttt{/\allowbreak{}token/\allowbreak{}verify/\allowbreak{}\{token\}} \\ \hline
\texttt{get-\allowbreak{}public-\allowbreak{}key} & GET & \texttt{/\allowbreak{}public-\allowbreak{}key} \\ \hline
\texttt{register-\allowbreak{}encryption-\allowbreak{}key} & POST & \texttt{/\allowbreak{}encryption-\allowbreak{}key} \\ \hline
\texttt{unregister-\allowbreak{}encryption-\allowbreak{}key} & DELETE & \texttt{/\allowbreak{}encryption-\allowbreak{}key} \\ \hline
\end{longtable}

\newpage

\section{Service Operations}
\label{sec:operations}

The interfaces of the \texttt{authorizationToken} service, its operations, data models and implementation are provided below.
In particular, each subsection first names the HTTP method and path used to call the operation, after which it names the abstract operation from the \texttt{authorizationToken} SD document, as well as input and output types.

All operations require the requester's identity, presented according to the provider's authentication policy: the client certificate (mutual TLS) and, with the \texttt{outsourced} policy, \texttt{Authorization: Bearer IDENTITY-TOKEN//<token>}; systems holding the Authentication secret key send \texttt{Bearer AUTHENTICATOR-KEY//<system>//<hash>}, and with the \texttt{declared} policy \texttt{Bearer SYSTEM//<system>} is sent. Request and response bodies are \texttt{application/json}.
Errors are answered with an \mref{ErrorMessage} body and the status codes \texttt{400 Bad Request} (invalid parameter), \texttt{401 Unauthorized} (missing or invalid identity), \texttt{403 Forbidden} (not permitted, e.g. by the management or blacklist filter), \texttt{404 Not Found}, \texttt{408 Request Timeout}, \texttt{423 Locked}, \texttt{500 Internal Server Error} and \texttt{503 Service Unavailable} (a system it depends on failed).

\osubsection{generate}{POST}{/\allowbreak{}consumerauthorization/\allowbreak{}authorization-token/\allowbreak{}generate}{\mref[Authorization\allowbreak{}Token\allowbreak{}Generation\allowbreak{}Request]{AuthorizationTokenGenerationRequest}}{\mref[Authorization\allowbreak{}Token\allowbreak{}Generation\allowbreak{}Response]{AuthorizationTokenGenerationResponse}}

Generates a token of the requested variant for the requester (as consumer) to the provider's target and operation, if authorized.

A successful call returns \texttt{201 Created} with the output in the body.

\begin{lstlisting}[label={lst:generate},caption={An \oref{generate} invocation.}]
POST /consumerauthorization/authorization-token/generate HTTP/1.1
Host: consumerauthorization:8445
Authorization: Bearer IDENTITY-TOKEN//<token>
Content-Type: application/json

{
  "tokenVariant": "TIME_LIMITED_TOKEN_AUTH",
  "provider": "TemperatureSensor",
  "targetType": "SERVICE_DEF",
  "target": "temperatureReading",
  "scope": "read"
}
\end{lstlisting}

\begin{lstlisting}[label={lst:generate-resp},caption={An \oref{generate} response.}]
{
  "tokenType": "TIME_LIMITED_TOKEN",
  "targetType": "SERVICE_DEF",
  "token": "<token>",
  "usageLimit": 5,
  "expiresAt": "2026-10-22T12:00:00Z"
}
\end{lstlisting}

\osubsection{verify}{GET}{/\allowbreak{}consumerauthorization/\allowbreak{}authorization-token/\allowbreak{}token/\allowbreak{}verify/\allowbreak{}\fparam{\{token\}}}{\pref{String}}{\mref[Authorization\allowbreak{}Token\allowbreak{}Verify\allowbreak{}Response]{AuthorizationTokenVerifyResponse}}

Verifies a token presented to the requester (as provider), counts a use of usage-limited tokens and returns its consumer, target and scope.

Path parameter: \fparam{token}.

A successful call returns \texttt{200 OK} with the output in the body.

\begin{lstlisting}[label={lst:verify},caption={An \oref{verify} invocation.}]
GET /consumerauthorization/authorization-token/token/verify/<token> HTTP/1.1
Host: consumerauthorization:8445
Authorization: Bearer IDENTITY-TOKEN//<token>
\end{lstlisting}

\begin{lstlisting}[label={lst:verify-resp},caption={An \oref{verify} response.}]
{
  "verified": true,
  "consumerCloud": "LOCAL",
  "consumer": "TemperatureDisplay",
  "targetType": "SERVICE_DEF",
  "target": "temperatureReading",
  "scope": "read"
}
\end{lstlisting}

\osubsection{get-public-key}{GET}{/\allowbreak{}consumerauthorization/\allowbreak{}authorization-token/\allowbreak{}public-key}{}{\pref{String}}

Returns the public key (Base64 X.509 Subject\allowbreak{}Public\allowbreak{}Key\allowbreak{}Info) with which providers verify JSON Web Tokens.

A successful call returns \texttt{200 OK} with the output in the body.

\begin{lstlisting}[label={lst:get-public-key},caption={An \oref{get-public-key} invocation.}]
GET /consumerauthorization/authorization-token/public-key HTTP/1.1
Host: consumerauthorization:8445
Authorization: Bearer IDENTITY-TOKEN//<token>
\end{lstlisting}

\begin{lstlisting}[label={lst:get-public-key-resp},caption={An \oref{get-public-key} response.}]
"string"
\end{lstlisting}

\osubsection{register-encryption-key}{POST}{/\allowbreak{}consumerauthorization/\allowbreak{}authorization-token/\allowbreak{}encryption-key}{\mref[Authorization\allowbreak{}Encryption\allowbreak{}Key\allowbreak{}Registration\allowbreak{}Request]{AuthorizationEncryptionKeyRegistrationRequest}}{\pref{String}}

Registers or replaces the requester's AES key for encrypting its self-contained tokens. Returns the initialization vector (CBC) or an empty string; 201 for a new key.

A successful call returns \texttt{201 Created} with the output in the body.

\begin{lstlisting}[label={lst:register-encryption-key},caption={An \oref{register-encryption-key} invocation.}]
POST /consumerauthorization/authorization-token/encryption-key HTTP/1.1
Host: consumerauthorization:8445
Authorization: Bearer IDENTITY-TOKEN//<token>
Content-Type: application/json

{
  "key": "0123456789abcdef",
  "algorithm": "AES/ECB/PKCS5Padding"
}
\end{lstlisting}

\begin{lstlisting}[label={lst:register-encryption-key-resp},caption={An \oref{register-encryption-key} response.}]
"string"
\end{lstlisting}

\osubsection{unregister-encryption-key}{DELETE}{/\allowbreak{}consumerauthorization/\allowbreak{}authorization-token/\allowbreak{}encryption-key}{}{}

Removes the requester's encryption key. Returns 200 if removed, 204 if there was none.

A successful call returns \texttt{200 OK} if the entry was removed, \texttt{204 No Content} if there was nothing to remove.

\begin{lstlisting}[label={lst:unregister-encryption-key},caption={An \oref{unregister-encryption-key} invocation.}]
DELETE /consumerauthorization/authorization-token/encryption-key HTTP/1.1
Host: consumerauthorization:8445
Authorization: Bearer IDENTITY-TOKEN//<token>
\end{lstlisting}

\newpage

\section{Data Models}
\label{sec:model}

Here, all data objects that can be part of the service calls associated with this service are listed in alphabetic order.
Note that each subsection, which describes one type of object, begins with the \textit{struct} keyword, which is meant to denote a JSON \pref{Object} that must contain certain fields, or names, with values conforming to explicitly named types.
Fields that are not needed may be omitted or \texttt{null} in requests; empty fields are omitted in responses.
As a complement to the primary types defined in this section, there is also a list of secondary types in Section \ref{sec:model:primitives}, which are used to represent things like hashes, identifiers and texts.

\msubsection{struct}{AuthorizationEncryptionKeyRegistrationRequest}

Registers the requester's key for encrypting its self-contained tokens.


\noindent\begin{tabularx}{\textwidth}{| >{\raggedright\arraybackslash}p{3.9cm} | >{\raggedright\arraybackslash\small}p{5.3cm} | X |} \hline
\rowcolor{gray!33} Object Field & Value Type & Description \\ \hline
"key" & \pref{String} & AES key of the provider (at least 16 bytes). \\ \hline
"algorithm" & \pref{String} & Encryption algorithm of the key, e.g. \texttt{AES/\allowbreak{}ECB/\allowbreak{}PKCS5Padding} or \texttt{AES/\allowbreak{}CBC/\allowbreak{}PKCS5Padding}. \\ \hline
\end{tabularx}



\msubsection{struct}{AuthorizationTokenGenerationRequest}

Request of the consumer for a token to one provider's service operation.


\noindent\begin{tabularx}{\textwidth}{| >{\raggedright\arraybackslash}p{3.9cm} | >{\raggedright\arraybackslash\small}p{5.3cm} | X |} \hline
\rowcolor{gray!33} Object Field & Value Type & Description \\ \hline
"token\allowbreak{}Variant" & \pref{String} & Requested token variant, as interface security policy (e.g. \texttt{TIME\_\allowbreak{}LIMITED\_\allowbreak{}TOKEN\_\allowbreak{}AUTH}, \texttt{RSA\_\allowbreak{}SHA512\_\allowbreak{}JSON\_\allowbreak{}WEB\_\allowbreak{}TOKEN\_\allowbreak{}AUTH}). \\ \hline
"provider" & \pref{String} & Provider system the token is for. \\ \hline
"target\allowbreak{}Type" & \pref{String} & Target type: \texttt{SERVICE\_\allowbreak{}DEF} or \texttt{EVENT\_\allowbreak{}TYPE}. \\ \hline
"target" & \pref{String} & Service definition (or event type). \\ \hline
"scope" & \pref{String} & Operation the token is for. \\ \hline
\end{tabularx}



\msubsection{struct}{AuthorizationTokenGenerationResponse}

A generated token.


\noindent\begin{tabularx}{\textwidth}{| >{\raggedright\arraybackslash}p{3.9cm} | >{\raggedright\arraybackslash\small}p{5.3cm} | X |} \hline
\rowcolor{gray!33} Object Field & Value Type & Description \\ \hline
"token\allowbreak{}Type" & \pref{AuthorizationTokenType} & Token type: \texttt{TIME\_\allowbreak{}LIMITED\_\allowbreak{}TOKEN}, \texttt{USAGE\_\allowbreak{}LIMITED\_\allowbreak{}TOKEN}, \texttt{SELF\_\allowbreak{}CONTAINED\_\allowbreak{}TOKEN} or \texttt{TRANSLATION\_\allowbreak{}BRIDGE\_\allowbreak{}TOKEN}. \\ \hline
"target\allowbreak{}Type" & \pref{AuthorizationTargetType} & Target type: \texttt{SERVICE\_\allowbreak{}DEF} or \texttt{EVENT\_\allowbreak{}TYPE}. \\ \hline
"token" & \pref{String} & The token (opaque string, Base64 or JSON Web Token). \\ \hline
"usage\allowbreak{}Limit" & \pref{Number} & Number of allowed uses (usage-limited tokens). \\ \hline
"expires\allowbreak{}At" & \pref{DateTime} & UTC expiry of time-limited tokens. \\ \hline
\end{tabularx}



\msubsection{struct}{AuthorizationTokenVerifyResponse}

Result of verifying a token presented to a provider.


\noindent\begin{tabularx}{\textwidth}{| >{\raggedright\arraybackslash}p{3.9cm} | >{\raggedright\arraybackslash\small}p{5.3cm} | X |} \hline
\rowcolor{gray!33} Object Field & Value Type & Description \\ \hline
"verified" & \pref{Boolean} & True when the token is valid. \\ \hline
"consumer\allowbreak{}Cloud" & \pref{String} & Cloud identifier of the consumer system. \\ \hline
"consumer" & \pref{String} & Name of the consumer system. \\ \hline
"target\allowbreak{}Type" & \pref{AuthorizationTargetType} & Target type: \texttt{SERVICE\_\allowbreak{}DEF} or \texttt{EVENT\_\allowbreak{}TYPE}. \\ \hline
"target" & \pref{String} & Target of the policy or token: a service definition or event type. \\ \hline
"scope" & \pref{String} & Operation the token or policy is limited to; empty for all operations. \\ \hline
\end{tabularx}



\msubsection{struct}{ErrorMessage}




\noindent\begin{tabularx}{\textwidth}{| >{\raggedright\arraybackslash}p{3.9cm} | >{\raggedright\arraybackslash\small}p{5.3cm} | X |} \hline
\rowcolor{gray!33} Object Field & Value Type & Description \\ \hline
"error\allowbreak{}Message" & \pref{String} &  \\ \hline
"error\allowbreak{}Code" & \pref{Number} &  \\ \hline
"exception\allowbreak{}Type" & \pref{ExceptionType} &  \\ \hline
"origin" & \pref{String} &  \\ \hline
\end{tabularx}



\subsection{Primitives}
\label{sec:model:primitives}

As all messages are encoded using the JSON format \cite{rfc8259}, the following primitive constructs, part of that standard, become available.
Note that the official standard is defined in terms of parsing rules, while this list only concerns syntactic information.
Furthermore, the \pref{Object} and \pref{Array} types are given optional generic type parameters, which are used in this document to signify when pair values or elements are expected to conform to certain types.

\noindent\begin{tabularx}{\textwidth}{| p{4.6cm} | X |} \hline
\rowcolor{gray!33} Type & Description \\ \hline
\pdef{AuthorizationTargetType} & One of \texttt{SERVICE\_\allowbreak{}DEF}, \texttt{EVENT\_\allowbreak{}TYPE} (JSON String). \\ \hline
\pdef{AuthorizationTokenType} & One of \texttt{TIME\_\allowbreak{}LIMITED\_\allowbreak{}TOKEN}, \texttt{USAGE\_\allowbreak{}LIMITED\_\allowbreak{}TOKEN}, \texttt{SELF\_\allowbreak{}CONTAINED\_\allowbreak{}TOKEN}, \texttt{TRANSLATION\_\allowbreak{}BRIDGE\_\allowbreak{}TOKEN} (JSON String). \\ \hline
\pdef{Boolean} & JSON Boolean, \texttt{true} or \texttt{false}. \\ \hline
\pdef{DateTime} & JSON String with an ISO 8601 UTC timestamp, e.g. \texttt{2026-10-22T12:00:00Z}. \\ \hline
\pdef{ExceptionType} & One of \texttt{ARROWHEAD}, \texttt{INVALID\_\allowbreak{}PARAMETER}, \texttt{AUTH}, \texttt{FORBIDDEN}, \texttt{DATA\_\allowbreak{}NOT\_\allowbreak{}FOUND}, \texttt{TIMEOUT}, \texttt{LOCKED}, \texttt{INTERNAL\_\allowbreak{}SERVER\_\allowbreak{}ERROR}, \texttt{EXTERNAL\_\allowbreak{}SERVER\_\allowbreak{}ERROR} (JSON String). \\ \hline
\pdef{Number} & JSON Number, an integer. \\ \hline
\pdef{Object} & JSON Object; \pref{Object}$<$A$>$ has string keys and values of type A. \\ \hline
\pdef{String} & JSON String. \\ \hline
\end{tabularx}

\newpage

\bibliographystyle{IEEEtran}
\bibliography{bibliography}

\newpage

\section{Revision History}
\subsection{Amendments}

\noindent\begin{tabularx}{\textwidth}{| p{1cm} | p{3cm} | p{2cm} | X | p{4cm} |} \hline
\rowcolor{gray!33} No. & Date & Version & Subject of Amendments & Author \\ \hline

1 & 2026-10-22 & \arrowversion & Document created from the Arrowhead LaTeX templates for the 5.2.2 Java Spring implementation, including the local extensions of this workspace. & Jerker Delsing \\ \hline
\end{tabularx}

\subsection{Quality Assurance}

\noindent\begin{tabularx}{\textwidth}{| p{1cm} | p{3cm} | p{2cm} | X |} \hline
\rowcolor{gray!33} No. & Date & Version & Approved by \\ \hline

1 & & \arrowversion  &  \\ \hline

\end{tabularx}

\end{document}
